\par\SourceLeaf{418}%
\Needspace{5\baselineskip}
\begin{center}3.\end{center}
\noindent\hspace*{1.2em}Beziehen wir die Lage der Punkte in der Ebene auf zwei zueinander senkrechte Koordinatenachsen, denen die Richtungen $0$ und $90^{0}$ entsprechen, so stellt die erste Koordinate den in Richtung der ersten Achse gemessenen Abstand von der zweiten Achse dar, die zweite hingegen den in Richtung der zweiten Achse gemessenen Abstand von der ersten. Bezeichnen ferner $x,y$ die Koordinaten eines beliebigen Punkts der gekrümmten Linie, $s$ ihre Länge von einem beliebigen Anfangspunkt bis zu diesem Punkt, $\varphi$ die Richtung der Tangente an diesem Punkt und $r$ den Krümmungsradius, so gilt
\[
\begin{aligned}
dx&=\cos\varphi.ds\\
dy&=\sin\varphi.ds\\
r&=\frac{ds}{d\varphi}.
\end{aligned}
\]

Wird die Gestalt der gekrümmten Linie durch die Gleichung $V=0$ dargestellt, wobei $V$ eine Funktion von $x$ und $y$ ist, und setzt man
\[
dV=pdx+qdy,
\]
so gilt auf der gekrümmten Linie
\[
pdx+qdy=0
\]
und damit
\[
p\cos\varphi+q\sin\varphi=0,
\]
woraus folgt
\[
\operatorname{tang}\varphi=-\frac pq.
\]
Ferner gilt
\[
\cos\varphi.dp+\sin\varphi.dq-(p\sin\varphi-q\cos\varphi)d\varphi=0.
\]
Setzt man also nach einem bekannten Lehrsatz
\[
\begin{aligned}
dp&=Pdx+Qdy\\
dq&=Qdx+Rdy,
\end{aligned}
\]
